\documentclass[10pt,letterpaper]{article}

\usepackage[utf8]{inputenc}
\usepackage[T1]{fontenc}
\usepackage[margin=1in]{geometry}
\usepackage[hidelinks]{hyperref}
\usepackage{microtype}
\usepackage{xcolor}
\usepackage{listings}
\usepackage{tabularx}
\usepackage{booktabs}
\usepackage{authblk}
\usepackage{enumitem}
\usepackage[backend=biber,style=numeric,sorting=none,maxbibnames=99]{biblatex}
\addbibresource{refs.bib}

\lstset{
    basicstyle=\ttfamily\small,
    breaklines=true,
    columns=fullflexible,
    keepspaces=true,
    showstringspaces=false,
    frame=single,
    rulecolor=\color{gray!40},
    backgroundcolor=\color{gray!5},
}

\title{Live-Patching as Moving-Target Defense:\\
Continuous In-Kernel Code Re-Randomization\\
with Post-Supersession Honey-Trap Tripwires}

\author[1]{Archis Gore}
\affil[1]{Encrypted Execution \quad \texttt{me@archisgore.com}}

\date{May 2026}

\begin{document}
\maketitle

\begin{abstract}
We observe that Linux's \texttt{CONFIG\_LIVEPATCH} subsystem ---
originally engineered as a reboot-free transport for CVE fixes --- can
be re-purposed as a continuous moving-target-defense (MTD) mechanism
that periodically replaces a kernel function with a semantically
identical copy at a different kernel-virtual address. We further
observe that under livepatch's atomic-replace mode, the
now-superseded code remains mapped and executable in module memory,
forming a natural surface for kprobe-based tripwires: any control
transfer to a previous generation's now-stale function address is, by
construction, illegitimate, and can be logged with caller information.
We demonstrate both ideas in a working open-source prototype,
\emph{evasive-linux}, that streams five livepatch generations under a
QEMU-booted Linux 6.6 kernel built with every modern randomization
knob enabled (KASLR base and memory, \texttt{RANDSTRUCT\_FULL},
\texttt{RANDOMIZE\_KSTACK\_OFFSET}, slab freelist randomization). A
final attacker module deliberately dispatches a call to the first
generation's superseded address; the honey-trap installed by the
second generation fires and the kernel logs the intrusion with the
caller identified. To our knowledge this is the first publicly
documented use of the Linux livepatch subsystem as a non-CVE
moving-target transport, and the first construction that converts
post-supersession dead code into a tripwire. Adelie [ASPLOS 2022] is
the closest prior work; it covers Linux modules only, uses bespoke
remapping rather than livepatch, and does not include a deception
component. We discuss the achievable cadence floor, the
re-randomization-coverage question, and the engineering obstacles
that have plausibly kept this construction unexplored to date.
Source, dossiers, and reproducible CI at
\url{https://github.com/encrypted-execution/evasive-linux}.
\end{abstract}

\section{Introduction}

Kernel address-space layout randomization (KASLR) in mainline Linux
is base-only and fixed at boot. Once the kernel image is decompressed
and relocated, a single secret offset hides every function for the
lifetime of the system. Over the last decade a string of side-channel
attacks --- prefetch timing~\cite{hund_prefetch,gruss_prefetch}, Intel
TSX speculative execution~\cite{drk}, Meltdown, and most recently
EntryBleed~\cite{entrybleed} --- have shown that a determined
attacker can recover the KASLR base with practical effort, even from
unprivileged userspace. Once recovered, every subsequent attempt on
the same machine benefits from the leak: the base does not change
until reboot, and reboot in production is expensive.

A natural response is to randomize \emph{continuously}: re-place
kernel functions in memory on a schedule, so that a leaked layout
expires before an attacker can build a code-reuse chain against it.
This idea has a productive history in userspace ---
TASR~\cite{tasr}, Shuffler~\cite{shuffler}, Readactor~\cite{readactor},
RuntimeASLR~\cite{runtimeaslr}, Isomeron~\cite{isomeron},
CodeArmor~\cite{codearmor}, Selfrando~\cite{selfrando} --- but
strikingly little uptake in the Linux kernel. The closest kernel-side
prior work, Adelie~\cite{adelie} (ASPLOS 2022), restricts itself to
loadable modules and bypasses upstream Linux infrastructure entirely
in favor of a bespoke module-remapping scheme. Function-granular
KASLR~\cite{fgkaslr} has spent six years out-of-tree. PaX
RANDKSTACK~\cite{pax_randkstack} re-randomizes the kernel stack base
on every syscall but does not touch \texttt{.text}.

This paper makes two observations that, taken together, change the
engineering calculus for continuous in-kernel code rerandomization:

\begin{enumerate}[leftmargin=*]
\item Linux already ships the plumbing. The livepatch subsystem
\cite{livepatch_docs} provides per-function code redirection via
ftrace, a per-task consistency model that guarantees no thread is
mid-execution in a function being swapped \cite{poimboeuf_consistency},
and an ``atomic replace'' mode in which a new patch supersedes every
previous one in a single transition. This infrastructure was engineered
to deliver CVE fixes without reboots; the same primitives also serve
as a moving-target transport. We are not aware of any prior published
work that proposes this re-use.

\item Atomic replace leaves the superseded function bodies resident
and executable. Once a new generation is installed, the kernel's
ftrace handler routes all legitimate callers of, say,
\texttt{cmdline\_proc\_show} to the new replacement; the
\emph{previous} replacement's text is no longer reached by any
legitimate path. An attacker whose memory disclosure leaked the
previous generation's address now holds a pointer that, if used,
indicates malice. This is a natural surface for a kprobe-based
tripwire, in the spirit of Crane et al.'s position
\cite{boobytrap} that every defense should be paired with a trap.
\end{enumerate}

We contribute:

\begin{itemize}[leftmargin=*]
\item \textbf{The first working prototype} of continuous code
re-randomization using Linux livepatch as the transport. Five
generations of an out-of-tree livepatch module, generated from a
template, atomic-replace one another in sequence; each lands its
replacement at a different kernel-virtual address under unmodified
upstream kernel infrastructure.
\item \textbf{A novel honey-trap construction.} Each new generation
registers a kprobe at the previous generation's now-superseded
function entry, whose \texttt{pre\_handler} logs any control transfer
to that location with full caller information. We demonstrate the
trap firing end-to-end against a simulated stale-pointer attacker.
\item \textbf{An end-to-end QEMU-reproducible demonstration} with
CI-validated GitHub Actions workflow that builds the kernel, the
livepatch modules, the busybox initramfs, boots the image, and
asserts the trap log line on every push.
\item \textbf{An honest survey of the prior art} (\S\ref{sec:related})
across the userspace continuous-rerand line, the much sparser kernel
line, and the deception-defense subfield, identifying two specific
unexplored angles that this work addresses.
\end{itemize}

\paragraph{Scope and posture.} This is a research prototype, not
production hardening. We are transparent about the engineering
obstacles that remain (\S\ref{sec:limitations}): the convergence-cadence
floor inherent in livepatch's consistency model, the coverage gap
created by inlining and tail-call optimization, the consistency
requirements of \texttt{\_\_jump\_table} and \texttt{.altinstructions}
across rerandomizations, and the information leak posed by
\texttt{/proc/kallsyms} and friends. We position evasive-linux as
\emph{additive} to encrypted-linux~\cite{encrypted_linux}, the
sibling project that applies the Encrypted Execution
thesis~\cite{encrypted_execution} at build time: encrypted-linux
diversifies the ABI in space (one dialect per build),
evasive-linux diversifies code placement in time (one layout per
cycle).

\section{Background: Linux Livepatch in Detail}
\label{sec:background}

A livepatch is an ordinary kernel module marked with
\texttt{MODULE\_INFO(livepatch, "Y")}. It declares a
\texttt{struct klp\_patch} containing one or more
\texttt{struct klp\_object} (one per target binary --- the running
\texttt{vmlinux} or any patched module), each holding
\texttt{struct klp\_func} entries that name the function to redirect
(\texttt{old\_name}), the replacement (\texttt{new\_func}), and an
optional \texttt{old\_sympos} to disambiguate static symbols of the
same name. On module \texttt{init}, \texttt{klp\_enable\_patch()}
is called.

\paragraph{Redirection via ftrace.} Livepatch does not overwrite the
prologue of the old function with a jump. It leans entirely on
dynamic-ftrace: every function compiled with \texttt{-pg -mfentry}
has a NOP slot at function entry that ftrace can atomically rewrite
to a \texttt{call \_\_fentry\_\_} trampoline. To redirect
\texttt{foo}, livepatch registers a \texttt{klp\_ftrace\_handler}
against \texttt{foo}'s fentry site; on entry the handler rewrites
the saved instruction pointer in \texttt{pt\_regs} to point at
\texttt{new\_foo}. The CPU then ``returns'' from the trampoline into
the new function. Crucially, \emph{the old function body is left
intact in \texttt{.text}}; only the entry NOP becomes an ftrace call.
This residue is the foundation of our honey-trap construction
(\S\ref{sec:design}).

\paragraph{The unified consistency model.} Three positions historically
contested how to make redirection consistent across a multi-CPU,
preemptible kernel: kpatch (Red Hat, 2014) used
\texttt{stop\_machine} plus a per-CPU stack walk; kGraft (SUSE, 2014)
used a per-task \texttt{TIF\_PATCH\_PENDING} flag with switchover at
syscall boundaries; Poimboeuf's unified model
\cite{poimboeuf_consistency} (merged in 4.12) combines them. On
\texttt{klp\_enable\_patch}, every task is tagged
\texttt{TIF\_PATCH\_PENDING}. A periodic worker performs reliable
stack-tracing on each task; if no patched function appears on the
stack, the task transitions to \texttt{KLP\_PATCHED}. Stuck tasks
transition naturally on their next syscall return. When the last
task has converged, \texttt{klp\_complete\_transition()} runs and
the sysfs \texttt{enabled} flag becomes 1. This convergence is the
floor on rerandomization cadence (\S\ref{sec:limitations}).

\paragraph{Atomic replace.} A patch declared with
\texttt{.replace = true} (introduced 2018) logically supersedes every
previously-loaded livepatch in a single transition. After convergence,
the older patch modules can be \texttt{rmmod}'d. This is the
fundamental primitive evasive-linux uses: a new generation atomically
takes over, the previous generation transitions to disabled state,
its sysfs \texttt{enabled} file becomes 0, and its module is now a
candidate for unload. Its \texttt{.text} remains mapped until
\texttt{rmmod} actually runs.

\paragraph{Symbol resolution.} A livepatch module needs to reference
kernel symbols that are not exported, including static functions.
Standard module loading cannot do that. The livepatch toolchain
therefore emits special relocations in
\texttt{.klp.rela.\{objname\}.\{section\}} sections that
\texttt{klp\_resolve\_symbols()} \cite{livepatch_docs} fixes up at
patch-enable time using kallsyms. For our purposes, the function
\texttt{cmdline\_proc\_show} is in kallsyms with
\texttt{CONFIG\_KALLSYMS\_ALL=y} (it is a static function in
\texttt{fs/proc/cmdline.c}), and a single \texttt{klp\_func} entry
naming it suffices.

\section{Threat Model}
\label{sec:threat}

We assume an attacker with the following capabilities:
\begin{itemize}[leftmargin=*]
\item A kernel-memory \emph{disclosure} primitive sufficient to
recover one or more kernel-text addresses (e.g., via a side channel
of the kind that defeats stock KASLR
\cite{hund_prefetch,gruss_prefetch,drk,entrybleed}, or via a
controlled-read primitive in a vulnerable subsystem).
\item A \emph{control-flow hijack} primitive separate from the
disclosure: either a function-pointer overwrite, a kernel ROP/JOP
chain, or a stale function-pointer cached in long-lived kernel state.
\item Sufficient time between the disclosure and the hijack to plan
the attack offline.
\end{itemize}

We do \emph{not} assume the attacker has:
\begin{itemize}[leftmargin=*]
\item Arbitrary kernel write-where (which would let them disable the
livepatch subsystem itself).
\item Read access to \texttt{/proc/kallsyms}, \texttt{/proc/kcore},
perf, ftrace, or BTF (which would defeat the defense by leaking the
new layout directly). We presume \texttt{lockdown=confidentiality}
or equivalent.
\item The ability to install their own livepatch (which would let
them reverse the redirection).
\end{itemize}

The defense aims to invalidate the disclosure between the leak and
the hijack, and to detect the hijack attempt when it happens.

The construction does \emph{not} defend against:
\begin{itemize}[leftmargin=*]
\item Just-in-time code-reuse (JIT-ROP) \cite{jit_rop,snow_zombie}
when the rerandomization cadence exceeds the attacker's leak-and-build
time. Ahmed et al.\ \cite{ahmed_rerand} measured this empirically:
function-level rerand must complete within roughly 1.5--2.4 seconds
to stay ahead of practical JIT-ROP. Our cadence is on the slow edge
of this regime.
\item Data-only attacks. We move code, not data.
\item Pure side channels that observe behavior without referencing
specific code addresses.
\end{itemize}

\section{Design}
\label{sec:design}

\subsection{Continuous Re-Randomization via Livepatch}

Treat the kernel function $f$ as a target of repeated, semantically
identical replacement. Construct a sequence of livepatch modules
$P_1, P_2, \ldots, P_n$, each containing a copy $f_i$ of the original
\texttt{f} (same source, recompiled into a different module). At
times $t_1 < t_2 < \ldots$, load $P_i$ with \texttt{.replace = true}.
The kernel's livepatch transition installs $f_i$ as the new
redirection target; all subsequent calls to \texttt{f} route to $f_i$.
$f_{i-1}$ remains resident at its allocation address, no longer
reachable through any ftrace-installed handler.

The address of each $f_i$ is determined by the kernel's vmalloc-based
module allocator. With \texttt{CONFIG\_RANDOMIZE\_MEMORY=y} the
module region's base is randomized at boot; each subsequent
\texttt{module\_alloc} call advances within the region. Consecutive
$f_i$ thus land at distinct kernel-virtual addresses without
additional engineering. We confirm this empirically
(\S\ref{sec:eval}).

\subsection{Honey-Trap on the Superseded Function}

After loading $P_{i+1}$, immediately install a kprobe at the entry of
$f_i$. The kprobe's \texttt{pre\_handler} logs the trap with the
faulting instruction pointer (\texttt{p->addr}) and the caller
(\texttt{regs->ip}, dereferenced via \texttt{\%pS} for symbolic
formatting). We use kprobes rather than direct text-patching because:

\begin{itemize}[leftmargin=*]
\item \texttt{register\_kprobe()} handles the W$\oplus$X enforcement
internally via \texttt{text\_poke\_bp}. A module cannot ordinarily
write to its own \texttt{.text} after load; kprobes is one of the
few well-supported paths.
\item The kprobe \texttt{pre\_handler} is a clean callback. It runs
on every hit at the address. The handler can return zero (let the
original instruction execute) or non-zero (skip), giving us policy
flexibility we exploit in \S\ref{sec:limitations}.
\item Unregistration is also clean: \texttt{unregister\_kprobe()}
restores the original instruction byte and removes the trap. We do
this on $P_{i+1}$'s \texttt{module\_exit}.
\end{itemize}

A subtle point: the kprobe installed by $P_{i+1}$ at $f_i$'s entry is
owned by $P_{i+1}$, not by $P_i$. When $P_{i+1}$ is later superseded
by $P_{i+2}$, the kprobe at $f_i$ remains armed (because $P_{i+1}$ is
still loaded, only disabled as a livepatch). The trap is alive for
as long as the module that armed it is loaded. On final teardown,
we unload modules in reverse order ($P_n$ first), and each module's
exit unregisters its kprobe before it goes.

\subsection{Why This Is Different from Direct Text-Poking}

A more direct construction would overwrite the superseded function's
body with \texttt{INT3} (\texttt{0xCC}) or \texttt{UD2}
(\texttt{0x0F 0x0B}) bytes via \texttt{text\_poke\_bp}, then install
a custom \texttt{die\_notifier} to recognize hits in our poisoned
regions. We tried this path; the engineering tax is significant
because:

\begin{enumerate}[leftmargin=*]
\item \texttt{text\_poke}, \texttt{text\_poke\_bp}, and
\texttt{set\_memory\_rw} are no longer exported to modules in
mainline kernels. Workarounds (\texttt{kallsyms\_lookup\_name}
hacks, CR0.WP toggling) are fragile and dependent on the kernel's
self-protection posture.
\item The \texttt{die\_notifier} approach demands a discriminator
that recognizes our poisoned regions without false-positiving on
legitimate BUG() or WARN() callsites.
\item The kprobe path achieves the same observable property (trap
fires, intrusion logged) with kernel-blessed plumbing.
\end{enumerate}

The kprobe construction we ship is therefore the practical baseline.
The INT3-carpet variant and a synthetic-gadget-bait variant remain
as future work for environments where finer per-byte control is
desired (e.g., to confuse offline gadget-discovery tools like
ROPgadget).

\section{Implementation}
\label{sec:impl}

The prototype is open source at
\url{https://github.com/encrypted-execution/evasive-linux}, Apache
2.0 licensed. Total non-research code is approximately 600 lines of
shell and C plus a 70-line livepatch source template.

\subsection{Kernel Build}

Built against Linux 6.6.30 LTS. Configuration is the upstream
\texttt{tinyconfig} plus two fragments:

\begin{itemize}[leftmargin=*]
\item \texttt{kernel/livepatch.config} ---
\texttt{CONFIG\_LIVEPATCH=y}, \texttt{CONFIG\_DYNAMIC\_FTRACE\_WITH\_REGS=y},
\texttt{CONFIG\_UNWINDER\_ORC=y}, \texttt{CONFIG\_KALLSYMS\_ALL=y},
\texttt{CONFIG\_KPROBES=y}, \texttt{CONFIG\_OPTPROBES=y},
\texttt{CONFIG\_MAGIC\_SYSRQ=y}, \texttt{CONFIG\_SAMPLE\_LIVEPATCH=m}.
\item \texttt{kernel/randomize.config} --- every modern randomization
knob: \texttt{CONFIG\_RANDOMIZE\_BASE=y} (KASLR base),
\texttt{CONFIG\_RANDOMIZE\_MEMORY=y} (memory region bases),
\texttt{CONFIG\_VMAP\_STACK=y}, \texttt{CONFIG\_RANDOMIZE\_KSTACK\_OFFSET=y}
(per-syscall kernel stack offset, the
RANDKSTACK~\cite{pax_randkstack} equivalent),
\texttt{CONFIG\_STACKPROTECTOR\_STRONG=y},
\texttt{CONFIG\_SLAB\_FREELIST\_RANDOM=y},
\texttt{CONFIG\_SHUFFLE\_PAGE\_ALLOCATOR=y}, and full GCC-plugin
struct-layout randomization (\texttt{CONFIG\_GCC\_PLUGIN\_RANDSTRUCT=y},
\texttt{CONFIG\_RANDSTRUCT\_FULL=y}). A per-build randstruct seed is
generated from \texttt{/dev/urandom} before the kernel is compiled.
\end{itemize}

The build runs inside a single \texttt{ubuntu:24.04}-based Docker
container (\texttt{docker/Dockerfile.kbuild}) that stages the
kernel source, busybox, and \texttt{gcc-13-plugin-dev} (the only
non-trivial dependency, required for the randstruct GCC plugin).

\subsection{Livepatch Module Template}

\texttt{livepatches/lp\_template.c} is the source template:

\begin{lstlisting}[language=C]
static int evasive_lp_@NN@_cmdline_show(struct seq_file *m,
                                        void *v) {
    seq_printf(m, "evasive-linux lp_@NN@ - @MSG@\n");
    return 0;
}
static struct klp_func funcs[] = {
    { .old_name = "cmdline_proc_show",
      .new_func = evasive_lp_@NN@_cmdline_show, }, { }
};
static struct klp_object objs[] = {
    { .funcs = funcs }, { }
};
static struct klp_patch patch = {
    .mod = THIS_MODULE,
    .objs = objs,
    .replace = true,           /* atomic-replace */
};
static int evasive_trap_pre_handler(struct kprobe *p,
                                    struct pt_regs *regs) {
    pr_warn("evasive-linux HONEY-TRAP: superseded patch @NN@ hit "
            "at %px (caller=%pS)\n",
            p->addr, (void *)regs->ip);
    return 0;
}
static struct kprobe evasive_trap_kp = {
    .pre_handler = evasive_trap_pre_handler,
};
static int __init lp_@NN@_init(void) {
    int ret;
    pr_info("evasive-linux lp_@NN@: new_func at %px\n",
            evasive_lp_@NN@_cmdline_show);
    ret = klp_enable_patch(&patch);
    if (ret) return ret;
    if (poison_addr) {
        evasive_trap_kp.addr = (kprobe_opcode_t *)poison_addr;
        ret = register_kprobe(&evasive_trap_kp);
        if (ret == 0)
            pr_info("HONEY-TRAP armed at %px\n",
                    (void *)poison_addr);
    }
    return 0;
}
\end{lstlisting}

The build orchestrator
(\texttt{scripts/build-continuous-demo.sh}) materializes
\texttt{lp\_01.c} through \texttt{lp\_05.c} from this template by
substituting \texttt{@NN@} and \texttt{@MSG@}, then invokes
\texttt{make -C /opt/linux M=\$PWD modules} against the just-built
kernel tree. The five resulting \texttt{.ko} files are placed into a
busybox initramfs.

\subsection{Demo Init Script and Attacker}

The initramfs \texttt{/init} loads the five generations in sequence:

\begin{lstlisting}[language=bash]
PREV=0
for nn in 01 02 03 04 05; do
    insmod /lib/modules/lp_${nn}.ko poison_addr=0x${PREV}
    CUR=$(dmesg | grep "lp_${nn}: new_func"\
        | tail -1 | awk '{print $NF}')
    PREV=${CUR}
done
# Demo 03b: invoke the very first generation's stale address.
insmod /lib/modules/attacker.ko addr=0x${FIRST_ADDR}
\end{lstlisting}

\texttt{attacker.ko} is intentionally minimal: it accepts an
\texttt{addr} module parameter, spawns a kthread, and dispatches an
indirect call through that address. We use a kthread rather than
\texttt{module\_init} context so that the inevitable NULL dereference
inside the stale \texttt{cmdline\_proc\_show} (which expects a valid
\texttt{seq\_file *}) oopses the kthread without killing init.

\subsection{Reproducibility}

A GitHub Actions workflow (\texttt{.github/workflows/ci.yml}) runs
on every push:

\begin{itemize}[leftmargin=*]
\item Three demo jobs (\texttt{demo-01}, \texttt{demo-02},
\texttt{demo-03}) each build their respective kernel, livepatch
modules, and initramfs, then boot in QEMU.
\item Each job asserts the expected \texttt{dmesg} signature.
demo-03's strict assertion is:
\texttt{grep -q "evasive-linux HONEY-TRAP: superseded patch 02 hit"}.
\item A hygiene job runs shellcheck across all shell scripts and
verifies the research-dossier and config-fragment manifest.
\end{itemize}

Wall time on \texttt{ubuntu-latest} runners: 4--7 minutes per demo
job. Anyone can verify the result with one command:
\texttt{make demo-03}.

\section{Evaluation}
\label{sec:eval}

We make claims about \emph{feasibility} and \emph{behavioral
correctness}, not yet about microbenchmark overhead at high
rerandomization frequencies. Open question 1
(\S\ref{sec:limitations}) is whether the latency floor is
practically usable; we have not yet measured it directly.

\subsection{Five Distinct Placements}

Captured on a clean QEMU boot:

\begin{lstlisting}
[Gen 01] insmod lp_01.ko (no honey-trap on first generation)
         placement: 0xffffffffa0000023
[Gen 02] insmod lp_02.ko poison_addr=0xffffffffa0000023
         placement: 0xffffffffa0007023
         trap:      HONEY-TRAP armed at 0xffffffffa0000023
[Gen 03] insmod lp_03.ko poison_addr=0xffffffffa0007023
         placement: 0xffffffffa0010023
         trap:      HONEY-TRAP armed at 0xffffffffa0007023
[Gen 04] insmod lp_04.ko poison_addr=0xffffffffa0010023
         placement: 0xffffffffa0019023
         trap:      HONEY-TRAP armed at 0xffffffffa0010023
[Gen 05] insmod lp_05.ko poison_addr=0xffffffffa0019023
         placement: 0xffffffffa0020023
         trap:      HONEY-TRAP armed at 0xffffffffa0019023
\end{lstlisting}

Each generation lands at a distinct kernel-virtual address. The
inter-generation delta (0x6000--0x9000 bytes) is the module allocator's
natural padding plus per-module sections; no extra entropy is
required beyond the underlying KASLR.

\subsection{Honey-Trap Fires}

Continuing the same boot, the attacker module is loaded and dispatches
a call to the very first generation's address:

\begin{lstlisting}
Loading attacker.ko addr=0xffffffffa0000023
  attacker: dispatching call to 0xffffffffa0000023 ...
  evasive-linux HONEY-TRAP: superseded patch 02 hit at
        0xffffffffa0000027 (caller=evasive_lp_01_cmdline_show
                                  +0x5/0x18 [lp_01])
\end{lstlisting}

The trap message confirms the kprobe installed by generation 02 (at
the entry of generation 01's superseded function) fired on the
attacker's indirect call. The address reported by the trap
(\texttt{0xffffffffa0000027}) is offset 4 from the function entry
(\texttt{0xffffffffa0000023}); the offset corresponds to the
\texttt{endbr64} instruction inserted by the kernel's IBT-enabled
build, after which the kprobe is allowed to attach.

\subsection{Clean Disarm}

On demo teardown, modules are \texttt{rmmod}'d in reverse order. Each
unload prints its trap-disarm message:

\begin{lstlisting}
evasive-linux lp_05: HONEY-TRAP at 0xffffffffa0019027 disarmed
evasive-linux lp_04: HONEY-TRAP at 0xffffffffa0010027 disarmed
evasive-linux lp_03: HONEY-TRAP at 0xffffffffa0007027 disarmed
evasive-linux lp_02: HONEY-TRAP at 0xffffffffa0000027 disarmed
\end{lstlisting}

\subsection{KASLR-on-KASLR}

The kernel is built with the full randomization stack
(\texttt{RANDSTRUCT\_FULL}, \texttt{RANDOMIZE\_BASE},
\texttt{RANDOMIZE\_MEMORY}, \texttt{VMAP\_STACK},
\texttt{RANDOMIZE\_KSTACK\_OFFSET}). Across boots, the module region's
base shifts, so the absolute generation addresses are different from
the example above; the inter-generation deltas remain similar. The
livepatch-driven re-randomization stacks atop boot-time randomization
multiplicatively: an attacker who somehow recovers KASLR base still
faces fresh per-generation placement.

\section{Limitations and Open Questions}
\label{sec:limitations}

\subsection{Cadence Floor}

The unified consistency model \cite{poimboeuf_consistency} converges
``usually within a few seconds'' per upstream documentation. The
lower bound is governed by long-sleeping kernel threads
(\texttt{kthreadd}, idle workers) and IPI cost. Sysfs exposes a
\texttt{force} knob to break stuck transitions. Practical cadence is
plausibly 1--10 seconds per cycle; sub-second is probably out of
reach with the current consistency model. JIT-ROP measurements
\cite{ahmed_rerand} place the danger threshold at 1.5--2.4 seconds
for function-level rerand, so this defense is on the slow edge of
the regime. Block-granular re-randomization within livepatch (in the
spirit of Remix \cite{remix}) is a natural next step.

\subsection{Coverage Gap}

Not every kernel function is re-randomizable. Subtract:

\begin{itemize}[leftmargin=*]
\item Functions inlined at every callsite. Per-function patching
becomes a no-op; transitive-closure analysis is required (kpatch's
classic problem).
\item Tail-call-optimized callers. A tail call (\texttt{jmp foo})
skips the fentry NOP; livepatch's redirect doesn't fire.
\texttt{-fno-optimize-sibling-calls} mitigates at a performance
cost.
\item Functions on any \texttt{kthread}'s perpetual sleep stack
(scheduler core, RCU, IRQ entry). The consistency model refuses to
transition while these are mid-execution; they are effectively
unpatchable.
\item Functions with active \texttt{kprobe} or \texttt{eBPF} attachments.
Livepatch and kprobes are mutually exclusive on the same function via
\texttt{FTRACE\_OPS\_FL\_IPMODIFY}. A re-randomization cycle must
either evict probes and reinstate them, or skip those functions.
\end{itemize}

Our hypothesis: 60--80\% of \texttt{vmlinux} by function count is
re-randomizable in the limit; the remainder defines the static
attack surface that this defense does not move.

\subsection{Consistency of Patched-In Addresses}

A re-randomization cycle that moves a function's code also
invalidates any place its address is hard-coded:

\begin{itemize}[leftmargin=*]
\item \texttt{\_\_jump\_table} entries (static keys / jump labels).
\item \texttt{.altinstructions} (CPU-feature-conditional code
patches).
\item Paravirt op tables (\texttt{pv\_ops}, inline-patched in
some configurations).
\item Function-pointer fields in long-lived kernel data
structures (\texttt{file\_operations}, \texttt{net\_device\_ops}, the
syscall table, etc.).
\end{itemize}

A complete implementation must walk and rewrite all of these on every
cycle. The current prototype handles only the simple case (a single
non-aliased function); generalization is the central engineering
challenge.

\subsection{Information Leakage}

\texttt{/proc/kallsyms}, \texttt{/proc/kcore}, perf samples, ftrace
ring-buffer entries, and BTF all expose the new layout to anyone
with appropriate privileges. The defense presumes
\texttt{lockdown=confidentiality} or equivalent posture; root-or-
equivalent disclosure defeats it. Whether the disclosure primitives
that motivated continuous rerand (e.g., EntryBleed
\cite{entrybleed}) survive lockdown is itself an open question.

\subsection{Memory and Module-Slot Pressure}

Every cycle leaves a stale module loaded. Atomic-replace plus
periodic \texttt{rmmod} keeps the steady-state count bounded, but
each cycle pays for new module allocation, BTF generation, and
kallsyms churn. Long-running systems may exhaust module slot space
or fragment the module address region; mitigation likely requires a
custom slot recycler.

\section{Related Work}
\label{sec:related}

\subsection{Userspace Continuous Code Rerandomization}

A productive decade of work: TASR \cite{tasr} re-randomizes on
I/O-boundary events on the insight that disclosure becomes dangerous
only when its result crosses a trust boundary. Shuffler \cite{shuffler}
runs an asynchronous shuffle thread in-process and demonstrates
millisecond-cadence on x86\_64 with 7--15\% overhead. Readactor
\cite{readactor} pairs randomization with execute-only memory and
trampoline indirection. Isomeron \cite{isomeron} flips a coin per
indirect call between two semantically identical replicas. CodeArmor
\cite{codearmor} virtualizes the entire code space and seeds honey
gadgets into it (the closest userspace cousin of our work). RuntimeASLR
\cite{runtimeaslr} remaps the post-\texttt{fork} child to defeat
BlindROP. Stabilizer \cite{stabilizer} introduced the rerand-as-PRNG
methodology for benchmarking. Smokestack \cite{smokestack} extends to
the stack. CoDaRR \cite{codarr} extends to data masks. None of these
target the kernel.

\subsection{Kernel-Side Rerandomization}

This is the sparse part of the literature. Giuffrida et al.\
\cite{giuffrida} demonstrate continuous rerand of the MINIX 3
microkernel with 5\% steady-state overhead, 10\% at 5-second cadence;
the microkernel design is essential to their approach. KASLR in
mainline Linux is base-only and broken by a sequence of side-channel
attacks \cite{hund_prefetch,gruss_prefetch,drk,entrybleed}. FGKASLR
\cite{fgkaslr} extends to function granularity at boot time only; it
has remained out-of-tree for six years. kMVX \cite{kmvx} runs
diversified kernels in parallel under N-variant execution to detect
leaks; kRX \cite{krx} provides execute-only memory; KHide \cite{khide}
adds disclosure-resistant diversification. All are orthogonal or
complementary to evasive-linux. PaX RANDKSTACK \cite{pax_randkstack}
shuffles the per-task kernel stack pointer every syscall; it is the
only continuous rerand mechanism shipping in production Linux and
serves as a useful precedent. Morpheus \cite{morpheus} demonstrates
50-millisecond pointer-encryption churn in a custom secure processor;
it survived a DARPA FETT bug bounty with $\sim$1\% overhead and is the
strongest existence proof that high-frequency MTD is affordable
\emph{in silicon}.

\paragraph{Adelie \cite{adelie}.} The closest published work. Adelie
re-randomizes Linux \emph{modules} via a bespoke remapping scheme
(custom PIC compilation plus an out-of-tree address-encryption layer)
on seconds-scale cadence. It does not target \texttt{vmlinux}, does
not use the livepatch subsystem, and does not include a deception
layer. We view evasive-linux as a complementary point in the design
space: where Adelie offers a self-contained module-only system,
evasive-linux re-uses Linux's own livepatch infrastructure, addresses
both modules and \texttt{vmlinux} in principle (limited in our
prototype to the simple case), and adds the honey-trap construction.

\subsection{Deception Defense}

Crane et al.'s position paper \cite{boobytrap} argues that every
defense should be paired with a trap: cancel the attack and log it.
HoneyGadget \cite{honeygadget} inserts gadget-shaped decoys at build
time and detects via Intel LBR sampling; CodeArmor \cite{codearmor}
seeds honey gadgets into its virtualized code space. Neither uses
the post-supersession dead code left by livepatch's atomic-replace as
the deception surface. Detection-only approaches (ROPscan, kBouncer
\cite{kbouncer}) verify return targets but do not trap. We are
unaware of prior published work that proposes the specific
construction: \emph{after a function is redirected by livepatch,
register a kprobe at the now-unreachable previous-generation entry as
an intrusion tripwire}.

\subsection{Build-Time Diversification}

Selfrando \cite{selfrando} performs load-time function permutation
in Tor Browser. encrypted-linux \cite{encrypted_linux}, the sibling
project, applies the Encrypted Execution thesis
\cite{encrypted_execution} at every contract point of the user-space
ABI and kernel UAPI --- syscall numbers, x86\_64 calling convention,
errno values, \texttt{/proc} schema, ELF \texttt{EI\_OSABI}, struct
layouts --- producing a different machine-level dialect per build.
That is the orthogonal axis: build-time spatial diversification.
evasive-linux adds runtime temporal diversification. Stacked, the two
defenses force an attacker to recover both the build's dialect and
the current generation's layout.

\section{Discussion}
\label{sec:discussion}

\paragraph{Why hasn't this been built?} Probably four reasons,
combined: (1) the \texttt{\_\_jump\_table} / \texttt{.altinstructions}
consistency engineering is non-trivial; (2) JIT-ROP \cite{jit_rop}
established a perception that any sub-instruction-granularity rerand
is insufficient against a determined attacker with disclosure
capability; (3) the Linux livepatch maintainer community has stayed
focused on the original CVE-fix use case; (4) Adelie \cite{adelie}
filled the seconds-cadence kernel-module niche just recently and
established a viable competitor that doesn't require livepatch
expertise to extend.

\paragraph{Why is the trap interesting on its own?} Because it
inverts a fundamental property of code-reuse attacks. In a static
KASLR world, a leaked function pointer is a useful pointer: it
identifies a callable entry. After livepatch supersession, a leaked
function pointer is an \emph{antagonistic signal}: its only legitimate
use vanished the moment the new generation took over. Any future
control flow toward that address is, by construction, malicious. The
defender's instrumentation cost is fixed (one kprobe per superseded
generation); the attacker's loss of optionality compounds with the
number of past generations.

\paragraph{Where the construction wants to go.} We see four natural
extensions: (a) drive the rerandomization loop from inside the kernel
(a workqueue that periodically inserts the next pre-built generation
via \texttt{request\_module}); (b) generate the next generation's
source code on demand at each tick, so the binary is genuinely
different not just located differently (composes cleanly with
encrypted-linux's per-build dialect); (c) add a synthetic-gadget
honey-bait variant that pollutes the offline ROP-discovery output of
tools like ROPgadget; (d) extend coverage by integrating Remix-style
\cite{remix} basic-block shuffling within each generation. Each is
an independent research project.

\section{Conclusion}
\label{sec:conclusion}

The Linux kernel ships, today, with the precise infrastructure
required to continuously re-randomize its own \texttt{.text}: per-function
ftrace-driven redirection, a transitional consistency model, and
atomic-replace for cumulative patches. It was engineered for CVE
hot-fixing, but it is also a moving-target transport. The dead code
left behind by atomic-replace is not waste; it is a perfectly-shaped
honey surface, and a kprobe at its entry converts a stale function
pointer in an attacker's working memory from a useful gadget into a
self-incriminating signal. Both observations appear to be unexplored
in the literature, and both fit naturally with mainline Linux.

We have demonstrated the construction end-to-end. The prototype is
small, open, CI-validated, and reproducible in fifteen minutes with
one command. The interesting questions --- how fast can we cycle,
how much of \texttt{vmlinux} can we cover, how to keep
\texttt{\_\_jump\_table} consistent across cycles --- are now
concrete engineering problems on top of a working baseline rather
than open theoretical questions.

\paragraph{Availability.} Source, research dossiers, and
reproducible CI:
\url{https://github.com/encrypted-execution/evasive-linux}.

\paragraph{Acknowledgments.} The author thanks the upstream livepatch
maintainers --- Josh Poimboeuf, Joe Lawrence, Jiri Kosina, Petr
Mladek, Miroslav Benes --- whose work made this construction possible.
The patent underlying this work (USPTO 10,733,303) is publicly
pledged to the public domain.

\printbibliography

\end{document}
